\chapter*{ЗАКЛЮЧЕНИЕ}
\addcontentsline{toc}{chapter}{ЗАКЛЮЧЕНИЕ}
\sloppy

Настоящая работа была посвящена разработке и апробации инструментов автоматизированной проверки заданий для курса по классическому машинному обучению, реализованного в НИУ ВШЭ с участием 200 студентов третьего курса.

В рамках работы были выполнены следующие задачи. Разработана рабочая программа курса, охватывающая восемь тематических разделов классического машинного обучения --- от линейных моделей и деревьев решений до ансамблевых методов и кластеризации --- и включающая систему оцениваемых заданий: домашние работы, активность на семинарах и семестровый проект. Реализовано расширение инструмента nbgrader, добавляющее новый тип ячейки LLM graded answer, который обеспечивает полную интеграцию БЯМ-оценщика в стандартный рабочий процесс преподавателя без изменения привычных команд. Разработана многоагентная система проверки проектных репозиториев, в которой каждый агент специализирован на отдельной категории рубрики, самостоятельно навигирует по репозиторию через ограниченный набор инструментов и изолирован от сбоев остальных агентов.

Проведённые эксперименты показали, что среди рассмотренных конфигураций наиболее сбалансированным по совокупности критериев --- accuracy, скорости и стоимости --- оказался оценщик на базе модели gpt-oss-120b без эталонного ответа в промпте: accuracy~0{,}93 при среднем времени отклика~5{,}09~с и суммарной стоимости~2{,}71~руб. за полный прогон датасета. Включение эталонного ответа в промпт не даёт статистически значимого прироста качества, но стабильно увеличивает стоимость и задержку.

